\section{Background}
\label{ch:background}

This chapter defines transition-graph terminology in Section~\ref{sec:terminology}.
Sections~\ref{sec:android} to~\ref{sec:dynamic-spa} review static extraction for Android,
navigation models for web applications and dynamic models for single-page applications (SPAs).
The chapter then covers source-level static analysis (Section~\ref{sec:asttools}), formal
foundations (Section~\ref{sec:formal}), Angular mechanisms (Section~\ref{sec:angular}) and
reactive programming's limits on static analysis (Section~\ref{sec:reactive}).
Section~\ref{sec:mutation} reviews mutation testing.
Section~\ref{sec:gap} summarises the chapter and identifies the research gap.

\subsection{User-interface transition graphs}
\label{sec:terminology}

A UI transition graph is a directed graph whose nodes represent user-visible application
states and whose edges represent transitions caused by user actions.
Figure~\ref{fig:utg-example} illustrates the model used here.

\begin{figure}[t]
\centering
\begin{tikzpicture}[node distance=14mm and 38mm]
  \node[screen] (home) {/};
  \node[screen, right=of home] (settings) {/settings};
  \node[screen, above=of home] (login) {/login};
  \node[screen, below=of home] (article) {/article/:slug};

  \draw[mustedge] (home) -- node[lbl, left] {\texttt{click} $\cdot$ must} (login);
  \draw[mayedge] (home.15) -- node[lbl, above=1mm] {\texttt{click} [\textit{requireAuth}] $\cdot$ may} (settings.165);
  \draw[mustedge] (settings.195) -- node[lbl, below=1mm] {\texttt{navigate} $\cdot$ must} (home.345);
  \draw[mayedge] (home) -- node[lbl, left] {\texttt{click} $\cdot$ may} (article);
\end{tikzpicture}
\caption{A small UI transition graph in the model of this thesis, with must and may edges.}
\label{fig:utg-example}
\par\smallskip{\footnotesize\noindent\emph{Note.} Each edge carries an event, an optional guard in brackets, and a modality: \emph{must} always fires when its code runs, \emph{may} is conditional (here on a route guard or a partly known target). In the measured graphs the unknown modality marks only edges into a sink for destinations that only the run time knows.}
\end{figure}


Research communities use different names for these graphs. In Android static analysis it is a
\emph{window transition graph} \cite{yang2015wtg}, \emph{activity transition graph} \cite{storydroid2019} or \emph{screen transition
graph} \cite{goalexplorer2019}. For Ajax (asynchronous JavaScript and XML) crawlers it is a \emph{state
flow graph} over DOM states \cite{crawljax2012}. In web
engineering it is a \emph{navigation model} or \emph{navigation state
machine} \cite{riccatonella2001,sitehopper2013}. This thesis uses \emph{UI transition graph}
for all of them, while recognising differences in state definitions and graph construction.

Two distinctions matter more than names. Extraction is either \emph{static}, from
source code or binaries without execution, or \emph{dynamic}, from observed runs. Edges
may or may not record \emph{guards}, the conditions under which transitions are available.
Table~\ref{tab:rw} classifies the reviewed work by both distinctions.

\begin{table}[tbp]\small
\caption{Prior work closest to this thesis, by target, extraction mode, guard representation and evaluation basis.}
\label{tab:rw}
\setlength{\tabcolsep}{4pt}
\begin{tabularx}{\textwidth}{@{}lYllY@{}}
\toprule
Work & Target & Static & Guards & Evaluation basis \\
\midrule
Yang et al.\ \cite{yang2015wtg,yang2018wtg} & Android windows & \checkmark & \xmark & 20 apps; manual on 6 \\
Yang et al.\ \cite{yang2015callbacks} & Android callbacks & \checkmark & context & manual vs.\ dynamic \\
StoryDroid \cite{storydroid2019} & Android activities & \checkmark & \xmark & hand-built, 10 apps \\
GoalExplorer \cite{goalexplorer2019} & Android screens & hybrid & \xmark & manual, 36 apps \\
ProMal \cite{promal2022} & Android windows & hybrid & \xmark & 40 apps, F1 score \\
UI-CTX \cite{uictx2025} & Android widgets & \checkmark & branches & 2000 labelled widgets \\
Ricca and Tonella \cite{riccatonella2001} & server-rendered pages & semi & \checkmark & manual \\
WebUml \cite{webuml2004} & server-rendered pages & hybrid & \checkmark & case study \\
Tkachuk and Rajan \cite{tkachuk2013patent} & Struts, JavaServer Pages (JSP) & \checkmark & outcome & not stated \\
Crawljax \cite{crawljax2012} & Ajax DOM states & \xmark & \xmark & manual \\
FeedEx, FRAGGEN \cite{feedex2013,fraggen2022} & DOM states & \xmark & \xmark & coverage \\
Kurant \cite{kurant2025} & React components & \checkmark & n/a & F1 score, no routing \\
Oshima and Kishi \cite{oshima2018model} & Angular pages & \checkmark & \xmark & 6 apps, faults \\
Rima et al.\ \cite{rima2026e2e} & Angular workflows & \checkmark & conditions & 7 apps, tests run \\
Zhang et al.\ \cite{zhang2025veriflow} & SPA routes, from bundles & \checkmark & \xmark & 4 apps (2 SPA), vs.\ Zed Attack Proxy (ZAP) \\
\bottomrule
\end{tabularx}
\par\smallskip{\footnotesize\noindent\emph{Note.} Evaluation basis is a reference, a metric or the tool compared with, as each work states it.}
\end{table}

\subsection{Static extraction for Android}
\label{sec:android}

Most static transition-graph research concerns Android. Yang et al.\ introduced the window
transition graph, with activities, dialogs and menus as window nodes and UI-thread
callback transitions as edges. The analysis models the active window stack, including back-button
behaviour \cite{yang2015wtg,yang2018wtg}. A related analysis builds a callback control-flow
graph over handler--widget pairs \cite{yang2015callbacks}.
StoryDroid builds an activity transition graph from \texttt{startActivity} calls \cite{storydroid2019}. GoalExplorer combines a statically derived screen
transition graph with dynamic exploration \cite{goalexplorer2019}. UI-CTX finds the code context of a widget's behaviour
with a backward data-flow analysis from the handler to the call that retrieves the widget \cite{uictx2025}.

These approaches attribute navigation to a screen through object resolution. A handler is
registered on a widget that belongs to exactly one window at run time, so resolving the
widget's identifier, usually by constant propagation of \texttt{getId()}, identifies that window
\cite{yang2015callbacks}. ProMal illustrates the limitation: imprecision in the reference
analysis aliases several widgets, assigning handlers to the wrong window and producing
infeasible transitions. Its authors combine static analysis with runtime traces and a
learned model \cite{promal2022}, using dynamic evidence rather than refining static attribution.

\subsection{Navigation models for web applications}
\label{sec:webmodels}

Web navigation models were first extracted from server-rendered sites. Ricca and Tonella
modelled pages, frames and forms connected by links and form submissions. Edges carry boolean
conditions on input values, derived semi-automatically using a crawler and analyst annotations
\cite{riccatonella2001}. WebUml combines static analysis of Hypertext Markup Language (HTML),
script and server-side pages with server runs to recover a state diagram. It derives guards
through control-flow analysis of server code
\cite{webuml2004}. A patent assigned to Fujitsu describes a fully static method
for Struts and JavaServer Pages (JSP) sources \cite{tkachuk2013patent}. SiteHopper builds abstracted
navigation state machines on the fly for model checking \cite{sitehopper2013}.

These approaches target server-rendered frameworks, where pages correspond to files and links
to anchor elements. They do not address client-side routers, component trees or dependency
injection (cf.~Section~\ref{sec:dynamic-spa}).

\subsection{Dynamic state models for single-page applications}
\label{sec:dynamic-spa}

With Ajax and client-side frameworks, model construction moved to run time. Crawljax executes
user events, compares the DOM before and after each event, and builds a state flow graph over
equivalence classes of DOM states \cite{crawljax2012}. Later work changed the exploration
strategy but not the model's source
\cite{feedex2013,fraggen2022}. The survey by Kertusha et al.\
reviews web testing research from 2014 to 2025 \cite{websurvey2025}. Every listed model-based
tool whose construction method was examined for this thesis executes the application to build
its model. The reasons are historical: Ajax applications lacked declarative route tables,
one uniform resource locator (URL) could represent many states, and source code did not
represent routes explicitly, although JavaScript static analysis existed \cite{jensen2009type}.

Client-side routing frameworks declare route tables as data, navigation as attributes or
method calls naming routes, and guards as named functions attached to routes. The literature
search found four static analyses of such frameworks: three analyse framework-specific source
code and one a compiled bundle. Kurant extracts a React component rendering graph but does not
address routing \cite{kurant2025}.

The other three address routing. Oshima and Kishi extract Angular page transitions and verify
them with the model checker SPIN (Simple Promela Interpreter) \cite{oshima2018model}. They analyse \texttt{RouterModule} declarations
and each component's navigable paths, resolving parent--child relationships among root, routed
and custom child components. Paths are attributed to the containing routed component, and the
root component's own paths to every page. The approach excludes \texttt{children} and route parameters
and does not treat data-binding changes as page transitions. The paper does not state how it
handles destinations that are nonliteral expressions, so no treatment is assumed here. Screens are
identified by routed-component names instead of addresses. Guards, lazy loading and named outlets
are not modelled. Navigation reached through a method call or dispatch is outside the
attribution mechanism. Without modalities, always-available and guarded transitions have
identical representations.

The evaluation uses four constructed applications and two published
applications of 1824 and 6893 lines, with six and twenty-three pages. It reports whether the
intended models were produced and hand-injected transition faults were detected in constructed
variants of the first sample. Losses on applications the extractor cannot fully analyse are
not measured, and measuring them is the subject of this thesis.

Two recent works statically analyse client-side router configuration for other purposes.
Rima et al.\ parse Angular sources and templates into a graph of navigation, interaction
sites and execution conditions, including guards. They compile workflows enumerated from
the graph into browser tests and use the workflow set as a coverage denominator
\cite{rima2026e2e}. They compare their output with a ground truth that a single auditor built
from the source and checked against their tool's artifacts.
The chapter defines a workflow as one trigger with its effects, from start routes to a single
terminal, and takes the active components of a route to be those activated by the route and its
ancestors, closed under composition. The further rules below are not in the chapter; they come
from the replication package's operational specification
\path{softscanner-sw/softscanner_fa}/\path{docs/paper/approach.md} and its code. Bootstrap
components are also active, and handlers are traced into same-class helpers to depth one and into
the first \texttt{subscribe} callback. The specification excludes router-outlet modelling and provider
resolution beyond module and root providers, and does not follow calls into injected services. The implementation nevertheless
traces calls into services whose class names match a router or navigation pattern, to depth one. Navigation inside these services is attributed to
their callers' screens.

Zhang et al.\ parse compiled JavaScript bundles, using declared routes as vertices and
navigation calls as edges. They check an access-control policy against screens reachable
from an authorised screen. Two of their four subjects are multi-page applications, analysed
by crawling instead. Their graph is an over-approximation and is not narrowed. The paper does
not explain source-screen attribution, and guards remain unmodelled runtime conditions
\cite{zhang2025veriflow}. Neither work classifies losses by construct
or measures them across a corpus. The problem is therefore neither new nor settled (cf.~Section~\ref{sec:gap}).

\subsection{Static analysis of JavaScript and TypeScript source}
\label{sec:asttools}

This thesis analyses Angular source with the TypeScript compiler
application programming interface (API), so it works on a source-level abstract syntax tree
rather than on a compiled artefact. This section reviews that style of tooling, the abstract
interpretation of JavaScript, what the declared types of a TypeScript program contribute, and
the evidence that frameworks limit such analyses more than the language does.

\subsubsection{Analyses over a source-level syntax tree}

Analysing the syntax tree of the source, instead of bytecode or an intermediate representation,
is a long-established style. Baxter et al.\ detect exact and near-miss code clones by comparing
abstract syntax trees with a variant of the common-subexpression detection used in compilers.
They apply the method to a C system of about 400\,000 lines \cite{baxter1998clone}. Working on
the tree keeps the analysis close to what the developer wrote, so each result can be reported at
a source location.

SPOON turns this style into a reusable library. It exposes a metamodel of Java source and lets
an analysis or a transformation be written in plain Java, without modifying a compiler or
learning a separate formalism \cite{pawlak2016spoon}. The TypeScript compiler API plays the same
role for TypeScript. It provides a typed syntax tree, a symbol table and a type checker over an
ordinary program. With them, an analysis can follow a constant reference in a route table to its
definition and read the declared type of an injected field. The cost of the style is that the
analysis sees the program as written, including the parts that a compiler later rewrites or
erases.

\subsubsection{Abstract interpretation of JavaScript}

Lattice-based analysis of JavaScript has a formal treatment. Jensen et al.\ present an analysis
based on abstract interpretation that infers detailed and sound type information for the full
language, including its object model and built-in functions. They propose it as a foundation for
developer tools \cite{jensen2009type}. Their work shows that the framework of
Section~\ref{sec:formal} applies to JavaScript at the level of individual values. That problem is
considerably harder than scope attribution, which uses the same machinery over a much coarser
domain of screens instead of abstractions of heap objects (cf.~Section~\ref{sec:scope}).

\subsubsection{What TypeScript's types contribute}

Angular applications are written in TypeScript, so an analysis can read declared types instead
of inferring them. The type system is deliberately unsound. Bierman et al.\ give a precise
definition of it \cite{bierman2014understanding}, and Richards et al.\ describe its annotations
as unsound by design \cite{richards2015concrete}. A declared type is therefore evidence about
the program and gives no guarantee about its execution.

Empirical studies point the same way. Gao et al.\ took bugs that had been fixed in the histories
of JavaScript projects, added type annotations to the buggy versions and checked whether a type
checker would have reported them. Flow and TypeScript each detected about 15\% of these bugs
\cite{gao2017type}. Bogner and Merkel mined 604 GitHub projects, 299 in JavaScript and 305 in
TypeScript. The TypeScript projects showed better code quality and understandability, but neither
lower bug proneness nor shorter bug resolution time. Heavier use of the \texttt{any} type
correlated weakly with worse measurements \cite{bogner2022type}. The declarations that describe
untyped libraries can themselves be wrong. TSCheck, which compares a declaration file with the
library it describes, found 142 errors in the declarations of 10 libraries
\cite{feldthaus2014checking}.

For an analysis of Angular code, a declared type narrows the set of candidate implementations of
a call, and a call graph can rely on that narrowing. The
type does not establish that the call reaches exactly one candidate. An unresolved or ambiguous
type must therefore widen the result instead of being trusted. In the graph model, this widening
is one source of the may tier.

\subsubsection{Frameworks as the practical obstacle}

In practice, the framework around the code limits a static analysis of a client-side application
more than the language does. Madsen et al.\ observe that JavaScript applications depend on large
frameworks and libraries whose behaviour an analysis must know. Hand-written models of that
behaviour are tedious to produce, incomplete and sometimes wrong. Their alternative infers library
behaviour from the way the application uses it, combining pointer analysis with such a \emph{use
analysis} \cite{madsen2013practical}.

A later analysis by Madsen et al.\ addresses implicit control flow directly. Their
\emph{event-based call graph} for Node.js represents the listeners that run when an event fires.
The analyses built on it find event-handling defects, such as events that no listener handles and
listeners registered too late \cite{madsen2015static}. An Angular application poses the same kind
of problem, with the router, template event bindings and observable subscriptions in place of the
Node.js event emitter.

Besides the navigation analyses of Oshima and Kishi (Section~\ref{sec:dynamic-spa}) and Rima et
al.\ \cite{rima2026e2e}, one static analysis of a framework in the Angular family is relevant
here. Aurebesh, by Ocariza et al., defines a consistency model for web applications built on a
model--view--controller framework \cite{ocariza2015detecting}. It detects identifier and type
inconsistencies between models, views and controllers, with a reported recall of 96.1\% and a
precision of 100\%, and it found 15 real bugs. Aurebesh targets AngularJS 1.x, whose template
language, component model and injector differ from those of the Angular versions analysed here.
The two frameworks share the problem of relating an identifier in a template to the code that
defines it. Aurebesh solves that problem for a consistency check, does not address navigation and
reports no transition graph.

\subsection{Formal foundations}
\label{sec:formal}

The analysis builds on three ideas from the literature: a transition system whose edges carry a
modality, a stated position on what soundness can mean for a practical analysis, and dataflow
analysis as the computational framework.

\subsubsection{Modal transition systems}

The \emph{modal transition system} of Larsen and Thomsen labels each transition as
either required, a \emph{must} edge, or merely permitted, a \emph{may} edge
\cite{larsen1988modal}. The graph model of this thesis (cf.~Section~\ref{sec:model}) borrows the names
\emph{must} and \emph{may} with a different meaning. There they are confidence labels on a
conditional runtime claim, not obligations on an implementation, and no refinement property of
modal transition systems is used. The model adds \emph{unknown} for edges whose modality could not
be determined.
The thesis claims soundness for the must tier alone, and the evaluation distinguishes the tiers
because may edges have weaker supporting evidence.

With one edge kind, an unresolved transition must either be included, reducing precision,
or omitted, reducing recall. Modalities record the evidence available, letting consumers
select a tier.

\subsubsection{Soundiness}
\label{sec:soundiness}

Livshits et al.\ argue that practical static analyses are inevitably sound only for a
language subset. They recommend declaring approximated features rather than claiming general
soundness \cite{livshits2015soundiness}.
Sections~\ref{sec:asttools} to~\ref{sec:angular} explain the limits: TypeScript annotations
are not guarantees, framework configuration determines call targets, and reactive firing
conditions are not written in the program text.

\subsubsection{Dataflow analysis and abstract interpretation}

Kildall expressed many program analyses as fixpoint computations over lattices of abstract
values. A join combines information, which propagates along graph edges until stable
\cite{kildall1973unified}. A finite-height lattice and monotone transfer functions ensure
termination and an order-independent result. Cousot and Cousot developed the general theory
of abstract interpretation \cite{cousot1977abstract}, and Nielson et al.\ provide a textbook
account \cite{nielson1999principles}. Scope attribution uses
this framework (cf.~Section~\ref{sec:scope}).

An instance needs a lattice of bounded height and a propagation graph. When unresolved calls leave
the propagation graph incomplete, the result omits information in a way the framework itself cannot
detect.

\subsection{Angular routing, dependency injection and reactive state}
\label{sec:angular}

The analysis must account for Angular-specific mechanisms,
described in the framework documentation \cite{angulardocs}. Listings~\ref{lst:routes}
to~\ref{lst:service} introduce a running example.

\begin{lstlisting}[language=TypeScript, caption={A route table with a guarded route, nested children, a lazy route table and a constant path.}, label={lst:routes}]
export const routes: Routes = [
  { path: '', component: HomeComponent },
  { path: 'settings', component: SettingsComponent,
    canActivate: [requireAuth],
    children: [{ path: 'profile', component: ProfileComponent }] },
  { path: 'admin', loadChildren: () => import('./admin/admin.routes') },
  { path: paths.about, component: AboutComponent },   // constant, not a literal
];
\end{lstlisting}

\begin{lstlisting}[language=AngularHTML, caption={The root template: the header sits next to the outlet.}, label={lst:shell}]
<app-header />                 <!-- shell: rendered on every screen -->
<router-outlet />              <!-- the active route renders here -->

<!-- header.component.html -->
<a routerLink="/settings">Settings</a>
@if (user()) { <a [routerLink]="['/profile', user().name]">Me</a> }
\end{lstlisting}
\par\smallskip{\footnotesize\noindent\emph{Note.} Its links are present on every screen but belong to none of them.}

\begin{lstlisting}[language=TypeScript, caption={Navigation from an injected service.}, label={lst:service}]
@Injectable({ providedIn: 'root' })
export class AuthService {
  private router = inject(Router);
  private session = inject(SessionStore);
  logout(): void {
    this.session.clear();
    this.router.navigate(['/login']);
  }
}
\end{lstlisting}
\par\smallskip{\footnotesize\noindent\emph{Note.} The call site belongs to no component, so its screens can only be found through the components that call \texttt{logout}.}


\subsubsection{Routes, components and navigation sites}

Angular declares screens through a tree of route objects. Each has a \texttt{path} pattern
and a \texttt{component}, \texttt{redirectTo} target or \texttt{children} list. A route with
children but no component provides a path prefix and no screen. Routes can load lazily through
\texttt{loadChildren} or \texttt{loadComponent}. Guards include \texttt{canMatch},
\texttt{canActivate}, \texttt{canActivateChild} and \texttt{canDeactivate}, and they return a
boolean, a redirect, or a Promise or Observable of either. When \texttt{canMatch} returns false,
the router tries the next matching route rather than cancelling navigation. Resolvers fetch
data that the router awaits before activation. Paths need not be string literals: resolving
\texttt{path: routes.about.path}, for example, requires tracing the constant to its definition.

Each screen is rendered by a component: a class and template. Templates instantiate other
components by selector and use \texttt{<router-outlet>} to mark where the router inserts the
active child component. Components above the outlet form the application shell, present on
every screen; those below a routed component are present only on that screen. Structural
directives such as \texttt{@if} and \texttt{*ngIf} make nesting conditional. Components can
also be created programmatically at run time.

A template expresses navigation with the directive \texttt{routerLink}, whereas class code
calls \texttt{Router.navigate} or \texttt{Router.navigateByUrl}. The option \texttt{relativeTo}
permits targets relative to the current route. Navigation sites can occur in components,
injectable services, guards, interceptors and resolvers. Their source screens are not
explicitly declared and must be derived.

\subsubsection{Dependency injection}
\label{sec:di}

Angular creates services and supplies them to the components and services that declare them as
dependencies. This mechanism is dependency injection, a form of inversion of control: a class does
not construct or look up its collaborators but receives them from a container
\cite{johnson1988reusable,fowler2004injection}. A dependency is requested by type or by an
injection token. The same token can be bound to different implementations in different injector
scopes, through \texttt{useClass}, \texttt{useExisting} or a factory function.

The pattern is widely used outside Angular. Yang, Tempero and Melton defined operational
criteria for dependency injection and detected its use in the source code of 34 open-source Java
applications \cite{yang2008empirical}. The configuration that the pattern introduces can itself be
analysed statically. Laigner et al.\ catalogue 12 dependency injection anti-patterns and detect
them with a static analyser at a relative recall of 92.19\%, and at least 9 of the anti-patterns
occur frequently in open-source and industrial projects \cite{laigner2022cataloging}. Dependency injection
configuration can therefore be read from the source, but it does not state which binding a given
call site receives.

For a static analysis, injection turns a simple method call into a dispatch problem. A call such
as \texttt{this.nav.goHome()} cannot be linked to the method it invokes until the type of
\texttt{nav} and the binding of that type in the relevant injector are resolved. Linking it is an
instance of the call graph construction problem known from object-oriented languages, where a
virtual call must be resolved to its possible targets. Class hierarchy analysis resolves such a
call from the declared type and its subclasses \cite{dean1995cha}. Rapid type analysis narrows
the result to the classes that the program instantiates \cite{bacon1996rta}, and points-to
analysis narrows it further by tracking which objects each variable may hold
\cite{andersen1994pointsto}. Grove and Chambers place these algorithms in a common framework that
trades precision against cost \cite{grove2001callgraph}. For JavaScript, Feldthaus et al.\ show
that an approximate call graph, which tracks functions by the property names they are stored
under, is precise enough for practical tools \cite{feldthaus2013jscg}.

Dependency injection adds a further difficulty, because the binding of a type to an
implementation is written in framework configuration instead of ordinary code. The Java
literature has measured what this costs a general analysis and what recovers the loss.
Table~\ref{tab:framework-models} lists the analyses that model framework behaviour explicitly.
Sridharan et al.\ report that the heavy use of reflection in frameworks makes static taint
analysis ineffective, and that adding support for one framework at a time does not scale. Their
tool, F4F, instead generates a specification of framework behaviour from the application code and
its configuration files. An unmodified taint engine consumes the specification, and the approach
found 525 new issues across nine benchmarks \cite{sridharan2011f4f}.

\begin{table}[tbp]\small
\caption{Static analyses that model framework behaviour explicitly, and the mechanism each of
them resolves.}
\label{tab:framework-models}
\begin{tabularx}{\textwidth}{@{}lYY@{}}
\toprule
Work & Framework & Mechanism modelled \\
\midrule
F4F \cite{sridharan2011f4f} & Java web frameworks & reflective calls, from a specification generated from code and configuration \\
\addlinespace
JackEE \cite{antoniadis2020enterprise} & Java enterprise frameworks & components and entry points, declaratively \\
\addlinespace
Jasmine \cite{chen2022jasmine} & Spring & dependency injection and aspect-oriented programming \\
\addlinespace
Aurebesh \cite{ocariza2015detecting} & AngularJS 1.x & identifiers shared between template and controller \\
\bottomrule
\end{tabularx}
\end{table}

Antoniadis et al.\ report that the major Java analysis frameworks offer almost no support for
enterprise applications. Their analysis, JackEE, captures the concepts that enterprise frameworks
share, such as components and entry points, in a declarative and largely framework-independent
form. Against the most precise conventional analysis, it measures about four times higher
completeness and about a sixfold speedup \cite{antoniadis2020enterprise}. Landman et al.\ set a
bound on what such modelling can achieve. In a literature review and corpus study of Java
reflection, they find that reflection cannot be ignored in 78\% of projects and that resolving it
requires unsound assumptions \cite{landman2017challenges}.

The closest analogue to resolving Angular's injectors is Jasmine, by Chen et al.\
\cite{chen2022jasmine}. Spring's core mechanisms, dependency injection and aspect-oriented
programming, rely on dynamic language features, so a conventional call graph of a Spring
application misses edges. Jasmine extends the Soot framework to make the call graph of a Spring
program more complete. On micro-benchmarks and 18 real Spring programs it improves on tools built
with Soot and with Doop. The structural problem is the one that Angular poses. An injected field has a declared type, the
binding of that type lives in framework configuration, and an analysis that does not read the
configuration cannot connect the call to a body.

Angular differs from Spring in where the configuration lives and how far its effect reaches.
Provider arrays belong to modules, routes or components, so the same token can resolve to
different implementations depending on where the injection happens. Resolution is relative to an
injector, which makes the provider hierarchy an input to the analysis.

\subsubsection{Reactive state}
\label{sec:signals}

Recent Angular versions include signals and effects \cite{angulareffects}. A signal is a value
that changes over time, and an effect is a computation that runs again whenever a signal it read
changes. Signals and effects are an instance of reactive programming, which goes back to
functional reactive programming \cite{elliott1997frp} and has been surveyed by Bainomugisha et
al.\ \cite{bainomugisha2013survey}. In a reactive system, the runtime records the dependencies
between values as they are read and propagates updates without explicit calls in the program
text.

A navigation call placed inside an effect therefore fires under a data dependence in the reactive
graph instead of a control-flow condition. Static analysis can identify the signals the effect
reads but cannot in general determine when the effect runs, and this thesis does not attempt to
bound it. Angular applications also use observables of Reactive Extensions for JavaScript (RxJS)
for the same purpose, in subscriptions to router events, in responses from the server and in
state containers (cf.~Section~\ref{sec:reactive}).

\subsection{Reactive programming and the limits of static dataflow}
\label{sec:reactive}

In reactive programming, the runtime decides when a computation runs. This section reviews the
propagation discipline that reactive runtimes implement and the evidence on how well developers
can reason about reactive code with tools.

\subsubsection{Propagation order and glitch freedom}

A reactive program is a graph of values in which an update to one value must propagate to the
values that depend on it. Cooper and Krishnamurthi embed dynamic dataflow in a call-by-value
language, so that an expression evaluates to a node in a dataflow graph
\cite{cooper2006embedding}. They name the central hazard. A \emph{glitch} occurs when a signal is
recomputed before all the signals it depends on are up to date, which exposes an intermediate
value that no consistent state of the program ever had. Their runtime prevents glitches by giving
each signal a height and updating signals in order of increasing height. The order is
topological, so the dependency graph must be acyclic.

Drechsler et al.\ carry glitch freedom into a distributed setting. Their change-propagation
algorithm needs no central knowledge of the dependency topology and avoids propagating updates
that cannot change a result \cite{drechsler2014distributed}. An Angular signal graph records its
dependencies as they are read (Section~\ref{sec:signals}), so it is the kind of structure over
which these propagation algorithms are defined.

For a static analysis, the order in which reactive computations run is a property of the graph
that the runtime builds, and that graph is built from the reads that actually happen. Static
analysis can over-approximate the set of signals an effect may read. The number of times the
effect runs and the application state at each run depend on which writes occur and in what
order, so they cannot in general be determined exactly from source. They could still be
over-approximated, but this analysis does not attempt it.

\subsubsection{Reactive code and the tools developers have}

Empirical studies separate reading reactive code from reasoning about its behaviour. Salvaneschi
et al.\ ran an experiment with 127 subjects that compared reactive code with object-oriented code
using the Observer pattern. Comprehension was significantly better for the reactive version
\cite{salvaneschi2017positive}. The same group identifies tool support as the weak side. A
conventional debugger shows a call stack, whereas the developer's question concerns a dataflow
graph, so reactive programs need debugging concepts of their own. Their Reactive Inspector
implements such concepts and outperformed conventional debugging in a controlled experiment
\cite{salvaneschi2016debugging}.

Studies of RxJS report the same gap. Banken et al.\ find that conventional debuggers offer no
support for reactive libraries such as RxJS, and that developers fall back on inserting console
logging. Their RxFiddle visualises the structure and the values of an RxJS dataflow, and in an
experiment with 111 developers it let them finish debugging tasks faster
\cite{banken2018debugging}. Alabor and Stolze interviewed 10 professional engineers and observed
4 of them debugging RxJS code. The engineers relied on manually added trace logging instead of
specialised tools. The study is a small workshop paper, so it suggests the practice without
establishing it \cite{alabor2020debugging}.

The surface that a tool would have to cover is large. Zimmerle et al.\ mined the use of the
reactive libraries RxJava, RxJS and RxSwift on GitHub and Stack Overflow. They found that 95.2\%
of the libraries' operators are used in practice \cite{zimmerle2022mining}.

Angular navigation analysis inherits this limit. An effect or subscription exposes a read set
in source, but that set does not exactly determine firing. The graph model therefore
specifies \emph{unknown} edges with the read set as symbolic guard, following the soundiness
position of Section~\ref{sec:formal} for unresolved constructs. The prototype emits no such
edges for effects or subscriptions. Its unknown edges have destinations it cannot resolve
to screens.

\subsection{Mutation testing}
\label{sec:mutation}

RQ4 asks whether the graph responds to changes in the source constructs it represents. Mutation
testing is applied here in an unusual direction. The mutants are inserted into the application
source, and the artefact whose response is measured is the extracted graph instead of a test
suite. This section reviews the technique and the evidence for its central assumption, its use on
JavaScript and web applications, and the reporting practice for such experiments.

\subsubsection{The coupling effect}

Mutation testing introduces small syntactic changes, called mutants, into a program and measures
the share that a test suite detects. It rests on the \emph{coupling effect} stated by DeMillo et
al.: tests that detect simple errors are often effective against more complex errors as well
\cite{demillo1978hints}. Without this assumption, a mutant would say nothing about the faults
that developers actually make.

Jia and Harman provide the standard survey of the field, covering approaches, tools and empirical
results. Their trend analysis suggests that the techniques and tools are reaching maturity
\cite{jia2011mutation}. Papadakis et al.\ survey the following decade of work on the fundamental
problems of the technique and list those that remain open \cite{papadakis2019mutation}. The
evidence that mutants are a usable proxy for real faults comes from Just et al. On 357 real
faults in 5 open-source applications, mutant detection
correlated significantly with real-fault detection, independently of code coverage
\cite{just2014mutants}.

\subsubsection{Mutation for JavaScript and web applications}

Mutation analysis has been adapted to JavaScript and to web applications, where the cost of
generating and running mutants is a central concern. Mirshokraie et al.\ use static and dynamic
analysis to direct mutant generation at code that is error-prone or likely to affect the output.
The guidance reduces the cost and the number of equivalent mutants, and the authors propose
mutation operators specific to JavaScript and the web \cite{mirshokraie2013efficient}. The journal
version adds a ranking of functions. In it, more than 93\% of the generated mutants are
non-equivalent, and more than 75\% of the surviving non-equivalent mutants lie in the top 30\% of
the ranked functions \cite{mirshokraie2015guided}.

Praphamontripong and Offutt target faults in the way web components interact with each other,
covering client-to-server and server-to-server interaction. They define mutation operators for
such faults and implement them in a tool called webMuJava. A small case study suggests that the
resulting tests find web application faults, and the work appeared at a workshop
\cite{praphamontripong2010applying}.

Stryker is a mutation testing framework for JavaScript and TypeScript, with further
implementations for C\# and Scala. Its documentation lists more than 30 supported mutations and a
design that works with any test runner \cite{stryker}. No peer-reviewed publication describing
Stryker was found, and the project cites none. Stryker counts a mutant as detected when a test
fails or times out with the mutant active \cite{stryker}. This thesis instead scores a mutant by
whether the extracted graph shows the change the mutation predicts.

\subsubsection{Reporting a mutation experiment}

Papadakis et al.\ also give best-practice advice on the use of mutation in empirical studies
\cite{papadakis2019mutation}. This thesis follows it and states the operators,
mutant generation and sampling procedures, exclusions and per-operator detection rates instead of
only an aggregate figure.

In each of the works above, the mutants evaluate a test suite for the application. The use of
mutation in this thesis belongs instead to change impact analysis. Arnold and Bohner compare impact analysis approaches by how they are used, how they work
internally and how effective they are \cite{arnold1993impact}, and the collection they edited is
the standard reference for the area \cite{bohner1996software}. Chianti is a classic tool
of this kind. It decomposes the difference between two versions of a program into atomic changes
and reports, for each affected test, the changes that affect it \cite{ren2004chianti}. Lehnert's
taxonomy classifies impact analysis approaches by finer criteria than the earlier frameworks
\cite{lehnert2011taxonomy}.

A transition graph whose edges carry source witnesses supports the same kind of query at the
level of screens. The mutation experiment checks whether the graph changes where its source
changes. It is a sanity check of the mapping from witnesses to source and does not measure
change-impact accuracy.

\subsection{Summary and the gap}
\label{sec:gap}

Static web navigation extraction has largely targeted server-rendered frameworks, where pages
are files and links are anchors. Client-side routing frameworks have mostly been modelled
dynamically. This literature search found three static extractions of client-side navigation
graphs. Two analyse framework-specific source code and one a compiled bundle, each for a
different consumer. Oshima and Kishi establish the problem addressed here, and
Section~\ref{sec:dynamic-spa} sets out what their approach leaves out. Zhang et al.\ use their graph without
narrowing it, as an over-approximated reachability relation for access-control checks
\cite{zhang2025veriflow}. Rima et al.\ use the enumerated workflow space as a generated test
suite's coverage denominator and leave out workflows that they prune \cite{rima2026e2e}.
Rima et al.'s tool already analyses nested route hierarchies, constrains edges with guards and
traces router-service calls to a bounded depth. No earlier work found in this review combines
cross-class and service attribution under provider bindings, which generalises that bounded
tracing, with modalities that separate unconditional transitions from guarded or scope-uncertain
ones, and with construct-level loss measurement on applications large enough to contain
difficult constructs.

Android provides the closest static techniques. They resolve navigation's owning object,
which usually identifies the window, and use runtime traces when resolution is imprecise.
Fragments, background tasks and aliased widgets are exceptions handled with dynamic evidence.
Web navigation sites often have no such object, which prevents direct reuse of Android
techniques.

Sections~\ref{sec:asttools} to~\ref{sec:mutation} identify established methods and remaining
limits. JavaScript abstract interpretation is established \cite{jensen2009type}, as is
source-level analysis over typed syntax trees \cite{baxter1998clone,pawlak2016spoon}.
Explicit framework models recover call edges missed by general analyses
\cite{sridharan2011f4f,antoniadis2020enterprise,chen2022jasmine}. Aurebesh addresses AngularJS
1.x template--controller consistency rather than navigation \cite{ocariza2015detecting}.
TypeScript types narrow dispatch targets without resolving them, because the type system is
deliberately unsound \cite{bierman2014understanding,richards2015concrete}. Reactive execution
depends on a runtime-built graph, limiting what source analysis determines
\cite{cooper2006embedding,banken2018debugging}. Mutation testing measures model responses to
source changes \cite{jia2011mutation,papadakis2019mutation}, and on web applications it has been
used to evaluate test suites \cite{mirshokraie2013efficient,praphamontripong2010applying}. For
static client-side navigation graphs, it has been applied only to hand-injected faults in
constructed applications \cite{oshima2018model}. This review found no operator-level mutation
study of such graphs.
